\section{Lecture 11: Solid Miscellany --- Structure, Homological Dimension, and the Constructible Topology}
\label{lec:11}

This lecture collects several complements to the solid theory before the
course moves on to a new topic. There are three parts. First, Clausen shows
that for a solid analytic ring cut out by a ring of integral elements
$(\ur A,A^{+})$ the derived category $D((\ur A,A^{+})^{\solid})$ really is the
derived category of its heart --- the derived $A^{+}$-solidification of the
compact projective generator lives in degree $0$ --- with everything reduced,
by a filtered-colimit argument, to the case of a ring finitely generated over
$\mathbb{Z}$. Second, for such a ring $R$ he analyses the abelian category
$\Solid_{R}$: it is $\Ind$ of its finitely presented objects, these form an
abelian (``coherent'') subcategory, the quasiseparated ones are exactly the
countable pro-systems of finitely generated discrete modules, and the compact
projective generator $\prod_{\mathbb{N}}R$ is flat. Third, when $R$ is
moreover regular of dimension $d$ he proves the expected homological bound
$\Ext^{i}=0$ for $i>d$ on finitely presented modules --- the one genuinely
delicate point being a depth/Auslander--Buchsbaum argument. The lecture closes
by pointing out a second, more radical way to localize $D((R,\mathbb{Z})^{
\solid})$: not over the valuative spectrum of \cref{lec:9} but over the
\emph{constructible} topology of $\Spec R$, producing idempotent algebras like
the higher local fields $\mathbb{Z}(\!(Y)\!)(\!(X)\!)$.

Throughout, ``ring'' means commutative ring, everything lives in light
condensed abelian groups, and we use freely the solid theory
$\Solid\subseteq\CondAbl$ (\cref{lec:5,lec:6}), the relative solid theory over
$\mathbb{Z}[T]$ (\cref{lec:7}), the abstract notion of analytic ring
(\cref{lec:8}) and its localization over the valuative spectrum (\cref{lec:9}).

\subsection{The solid analytic ring of a pair \texorpdfstring{$(\ur A,A^{+})$}{(A,A+)}, revisited}

Recall from \cref{lec:8,lec:9} the following construction. Let $\ur A$ be a
solid ring and $A^{+}\subseteq A^{\circ}\subseteq\ur A(\ast)$ a subring of
power-bounded elements. To the pair one attaches a solid analytic ring
$(\ur A,A^{+})^{\solid}$, specified on the abelian level by the full
subcategory
\begin{equation}
\label{eq:11:mod-pair}
\Mod_{(\ur A,A^{+})}\ \subseteq\ \Mod_{\ur A}(\Solid_{\mathbb{Z}})
\end{equation}
of those modules obtained by forcing every $f\in A^{+}$ to become a solidified
variable (i.e.\ requiring $\mathbb{Z}[T]$-solidity along $T\mapsto f$ for all
$f\in A^{+}$; \cref{def:9:solid-structure}). The hypothesis $A^{+}\subseteq
A^{\circ}$ is what makes $\ur A$ itself complete, so that $\ur A\in\Mod_{(\ur
A,A^{+})}$.

\begin{remark}[When $A^{+}$ is not inside a ring of definition]
\label{rmk:11:completion-aside}
If one drops the power-boundedness $A^{+}\subseteq A^{\circ}$, then enforcing
\eqref{eq:11:mod-pair} may throw $\ur A$ out of the category; one must then apply
the localization (solidification) and obtain a new underlying ring --- a
completion of $\ur A$ in the appropriate sense.
\end{remark}

On the derived level (\cref{def:8:derived-cat}) one has
\begin{equation}
\label{eq:11:derived-pair}
D\bigl((\ur A,A^{+})^{\solid}\bigr)\ \subseteq\ D\bigl(\Mod_{\ur A}(\Solid_{
\mathbb{Z}})\bigr),
\end{equation}
the full subcategory of objects whose homology lies in $\Mod_{(\ur A,A^{+})}$ in
every degree. The inclusion has a left adjoint --- the derived
$A^{+}$-solidification --- with the usual formal properties (symmetric
monoidal, colimit-preserving). The subtlety, flagged in \cref{lec:8}, is that
in the generality of an arbitrary analytic ring \eqref{eq:11:derived-pair} need
\emph{not} be the derived category of the abelian category
\eqref{eq:11:mod-pair}: the obstruction is that applying derived solidification
to a basic object may produce homology outside degree $0$.

The abelian category $\Mod_{(\ur A,A^{+})}$ has a compact projective generator,
namely
\begin{equation}
\label{eq:11:abelian-generator}
\Bigl(\,\textstyle\prod_{\mathbb{N}}\mathbb{Z}\,\Bigr)\otimes_{\mathbb{Z},\solid}
\ur A,
\end{equation}
the compact projective generator $\prod_{\mathbb{N}}\mathbb{Z}$ of
$\Solid_{\mathbb{Z}}$ (\cref{cor:5:P-solid}) base-changed to $\ur A$. Since
$\prod_{\mathbb{N}}\mathbb{Z}$ is flat in $\Solid_{\mathbb{Z}}$
(\cref{cor:6:flat}), it is immaterial whether one uses the ordinary or derived
solid tensor product here. A compact generator of the derived category
\eqref{eq:11:derived-pair} is obtained from \eqref{eq:11:abelian-generator} by
applying the derived $A^{+}$-solidification.

\begin{theorem}[The derived category is the derived category of the heart]
\label{thm:11:degree-zero}
For $(\ur A,A^{+})$ as above, the derived $A^{+}$-solidification of
$(\prod_{\mathbb{N}}\mathbb{Z})\otimes_{\mathbb{Z},\solid}\ur A$ lives in degree
$0$. Consequently it is a compact projective generator of the abelian category
$\Mod_{(\ur A,A^{+})^{\solid}}$, and
\begin{equation}
\label{eq:11:D-equals-Dheart}
D\bigl(\Mod_{(\ur A,A^{+})^{\solid}}\bigr)\ =\ D\bigl((\ur A,A^{+})^{\solid}
\bigr).
\end{equation}
\end{theorem}

Thus, as long as the underlying solid ring $\ur A$ is not itself derived, the
whole theory is ``flat'', i.e.\ determined at the abelian level.

\begin{remark}[This settles a point left open earlier]
\label{rmk:11:settled}
In \cref{lec:8} it was not clear that \eqref{eq:11:D-equals-Dheart} holds; that
is exactly what \cref{thm:11:degree-zero} establishes.
\end{remark}

\subsection{Reduction to finite type over \texorpdfstring{$\mathbb{Z}$}{Z}}

Being derived $A^{+}$-solid is an element-wise condition, so it can be tested
on the finitely generated subrings of $A^{+}$.

\begin{proposition}[Solidification through finite-type subrings]
\label{prop:11:filtered-colimit}
Derived $A^{+}$-solidity is equivalent to derived $R$-solidity for all
$R\subseteq A^{+}$ finite type over $\mathbb{Z}$, and for any $M$ the derived
$A^{+}$-solidification is the filtered colimit
\begin{equation}
\label{eq:11:filtered-colimit}
M^{\mathbb{L}A^{+}\text{-}\solid}\ =\ \varinjlim_{\substack{R\subseteq A^{+}\\
R\ \mathrm{f.t.}/\mathbb{Z}}} M^{\mathbb{L}R\text{-}\solid},
\end{equation}
with structure map induced from $M$ (regarding $M$ as an $R$-module via
$R\hookrightarrow A^{+}\to\ur A$).
\end{proposition}

\begin{proof}[Sketch]
The right-hand side of \eqref{eq:11:filtered-colimit} is derived
$A^{+}$-solid: for fixed $R$, every term past the $R$-term is even derived
$R'$-solid for some $R'\supseteq R$ (hence $R$-solid), so a cofinal subsystem
is $R$-solid, and a filtered colimit of $R$-solid objects is $R$-solid; thus
the colimit is $R$-solid for every $R$, hence $A^{+}$-solid. And mapping out to
any derived $A^{+}$-solid target does not distinguish $M$ from the colimit, by
the same $R$-by-$R$ argument.
\end{proof}

By \cref{prop:11:filtered-colimit}, \cref{thm:11:degree-zero} reduces to
showing, for $R$ finite type over $\mathbb{Z}$, that the derived
$R$-solidification of $(\prod_{\mathbb{N}}\mathbb{Z})\otimes_{\mathbb{Z},\solid}
\ur A$ lives in degree $0$. Doing this in two steps --- base change to
$\Solid_{R}$, then tensor up to $\ur A$ over $\Solid_{R}$ --- reduces further
to the two statements
\begin{equation}
\label{eq:11:two-statements}
\Bigl(\,\textstyle\prod_{\mathbb{N}}\mathbb{Z}\otimes_{\mathbb{Z},\solid}R\,
\Bigr)^{\mathbb{L}R\text{-}\solid}\ \text{lives in degree }0
\quad\text{and is flat for the }\Solid_{R}\text{-tensor product.}
\end{equation}

Before analysing \eqref{eq:11:two-statements}, we record a general
consequence of the filtered-colimit picture.

\begin{corollary}[All completion happens at finite type]
\label{cor:11:completion-finite-type}
For an arbitrary commutative ring $R$, write $\Solid_{R}:=\Mod_{(R,R)^{\solid}}$
(so that every element of $R$ is solidified). Then
\begin{equation}
\label{eq:11:ind-fp}
\Solid_{R}\ =\ \Ind\bigl(\Solid_{R}^{\mathrm{fp}}\bigr),
\end{equation}
where the finitely presented objects are the cokernels of maps between the
compact projective generator $\prod_{\mathbb{N}}R$. Moreover the whole
finitely presented theory reduces to the finite-type case: all of the
completion is happening at the level of finite-type subrings of $R$, and the
rest is an algebraic filtered colimit.
\end{corollary}

\subsection{The generator over a finite-type ring}

\begin{proposition}[The generator is a product of copies of $R$]
\label{prop:11:generator}
Let $R$ be finite type over $\mathbb{Z}$. Then
\begin{equation}
\label{eq:11:generator-product}
\Bigl(\,\textstyle\prod_{\mathbb{N}}\mathbb{Z}\otimes_{\mathbb{Z},\solid}R\,
\Bigr)^{\mathbb{L}R\text{-}\solid}\ =\ \textstyle\prod_{\mathbb{N}}R
\end{equation}
$($an unlabelled index set is always countable$)$. For a general $R$, not of
finite type, the answer is instead the smaller object
\begin{equation}
\label{eq:11:generator-union}
\Bigl(\,\textstyle\prod_{\mathbb{N}}\mathbb{Z}\otimes_{\mathbb{Z},\solid}R\,
\Bigr)^{\mathbb{L}R\text{-}\solid}\ =\ \bigcup_{\substack{R'\subseteq R\\ R'\
\mathrm{f.t.}}} \textstyle\prod_{\mathbb{N}}R'.
\end{equation}
\end{proposition}

\begin{proof}
First take $R=\mathbb{Z}[X_{1},\dots,X_{n}]$. Recall (\cref{prop:7:pushforward,cor:7:generator}) that the derived $\mathbb{Z}[T]$-solidification of a module
tensored with $\mathbb{Z}[T]$ commutes with all limits and colimits and sends
$\mathbb{Z}$ to $\mathbb{Z}[T]$; the various solidifications commute, and once
$M$ is $X_{i}$-solid it remains so after solidifying in the other variables.
Solidifying one variable at a time therefore throws each $X_{i}$ into the
product:
\begin{equation}
\label{eq:11:polynomial-solid}
\Bigl(\,\textstyle\prod_{\mathbb{N}}\mathbb{Z}\otimes_{\mathbb{Z},\solid}
\mathbb{Z}[X_{1},\dots,X_{n}]\Bigr)^{X_{1}\text{-}\solid}
=\textstyle\prod_{\mathbb{N}}\mathbb{Z}[X_{1}]\otimes_{\mathbb{Z}[X_{1}],\solid}
\mathbb{Z}[X_{2},\dots,X_{n}]=\cdots=\textstyle\prod_{\mathbb{N}}\mathbb{Z}
[X_{1},\dots,X_{n}].
\end{equation}
In general, choose a surjection $\mathbb{Z}[X_{1},\dots,X_{n}]\twoheadrightarrow
R$. Solidifying in $\mathbb{Z}$ and then base-changing algebraically (in the
derived sense) from $\mathbb{Z}[X_{1},\dots,X_{n}]$ to $R$, one has already
solidified with respect to a generating set of $R$, so the object is already
$R$-solid and it remains to compute
\begin{equation}
\label{eq:11:base-change-to-R}
\Bigl(\,\textstyle\prod_{\mathbb{N}}\mathbb{Z}[X_{1},\dots,X_{n}]\Bigr)\otimes^{
\mathbb{L}}_{\mathbb{Z}[X_{1},\dots,X_{n}]}R.
\end{equation}
But $\mathbb{Z}[X_{1},\dots,X_{n}]$ is noetherian, so $R$ admits a (possibly
infinite) resolution by finite free $\mathbb{Z}[X_{1},\dots,X_{n}]$-modules.
For a finite free module the product manifestly commutes with the base change;
passing to the class of $\mathbb{Z}[X_{1},\dots,X_{n}]$-modules for which the
claim holds --- which contains the ring and is closed under the relevant
colimits --- gives
\begin{equation}
\label{eq:11:resolve-R}
\Bigl(\,\textstyle\prod_{\mathbb{N}}\mathbb{Z}[X_{1},\dots,X_{n}]\Bigr)\otimes^{
\mathbb{L}}_{\mathbb{Z}[X_{1},\dots,X_{n}]}R\ \xrightarrow{\ \sim\ }\
\textstyle\prod_{\mathbb{N}}R,
\end{equation}
as an equivalence in the derived category, computed by resolving $R$ and
applying the colimit-preserving tensor functor. (The finiteness of the ideal
of relations is what turns the naive limit of quotients into an honest product
of copies of $R$.)
\end{proof}

\subsection{Structure of \texorpdfstring{$\Solid_{R}$}{Solid\_R} for \texorpdfstring{$R$}{R} finite type}

\begin{theorem}[Coherence]
\label{thm:11:coherence}
Let $R$ be finite type over $\mathbb{Z}$. Then:
\begin{enumerate}
\item $\Solid_{R}=\Ind(\Solid_{R}^{\mathrm{fp}})$, where the finitely
presented objects $\Solid_{R}^{\mathrm{fp}}$ are the cokernels of maps between
copies of the compact projective generator $\prod_{\mathbb{N}}R$;
\item $\Solid_{R}^{\mathrm{fp}}$ is an abelian subcategory of $\Solid_{R}$,
closed under kernels, cokernels and extensions.
\end{enumerate}
\end{theorem}

Part (1) is formal from the existence of a compact projective generator. Part
(2) is a coherence statement, and its proof mirrors classical commutative
algebra. Recall that a discrete ring is \emph{coherent} when its finitely
presented modules form an abelian subcategory, and that this reduces --- by
manipulating short exact sequences --- to the single point that every finitely
generated ideal is finitely presented. The same reduction here shows it
suffices to check that
\begin{equation}
\label{eq:11:coherence-reduction}
\text{every finitely generated subobject of }\textstyle\prod_{\mathbb{N}}R\
\text{is finitely presented.}
\end{equation}
Any subobject of $\prod_{\mathbb{N}}R$ is quasiseparated (Hausdorff), being a
subobject of a quasiseparated object, so \eqref{eq:11:coherence-reduction} in
turn follows from: every finitely generated \emph{quasiseparated} solid
$R$-module is finitely presented. This is the content of the next lemma, which
identifies the quasiseparated finitely presented modules with the classical
pro-systems.

\begin{lemma}[Quasiseparated finitely presented $=$ pro-(f.g.\ modules)]
\label{lem:11:qsep-pro}
Let $M\in\Solid_{R}$ be quasiseparated. Then the following are equivalent:
\begin{enumerate}
\item $M$ is finitely presented;
\item $M$ is finitely generated;
\item $M=\varprojlim_{n}M_{n}$, where each $M_{n}$ is a finitely generated
discrete $R$-module and the transition maps are surjective.
\end{enumerate}
Moreover the category of such $M$ is the countable pro-category
\begin{equation}
\label{eq:11:pro-category}
\mathrm{Pro}_{\mathbb{N}_{0}}\bigl(\{\text{finitely generated }R\text{-modules}
\}\bigr),
\end{equation}
i.e.\ morphisms between such inverse limits are the maps of pro-objects.
\end{lemma}

\begin{proof}
The nontrivial implications are (2)$\Rightarrow$(3) and (3)$\Rightarrow$(1).

For (2)$\Rightarrow$(3), a finite generation gives a surjection
$\prod_{\mathbb{N}}R\twoheadrightarrow M$ with kernel $K$:
\[
\begin{tikzcd}[column sep=large]
0 \arrow[r] & K \arrow[r, hook] & \prod_{\mathbb{N}}R \arrow[r, two heads] & M
\arrow[r] & 0.
\end{tikzcd}
\]
Since $M$ is quasiseparated, the inclusion $K\hookrightarrow\prod_{\mathbb{N}}R$
is a quasicompact map, so $K$ is quasiseparated: a filtered union of compact
Hausdorff subspaces, on each of which the inclusion is a closed embedding. As
$\prod_{\mathbb{N}}R$ is metrizable (sequential), being closed on each compact
subset is the same as being closed, so $K$ is just a closed submodule of
$\prod_{\mathbb{N}}R$ for the product topology. Projecting to the first $n$
coordinates and letting $K_{n}\subseteq\prod_{i\le n}R$ be the image, one gets
$K=\varprojlim_{n}K_{n}$ with surjective transition maps, and analogously
$M=\varprojlim_{n}M_{n}$ with $M_{n}=\prod_{i\le n}R/K_{n}$; both are of the
form (3).

For (3)$\Rightarrow$(1) it suffices to prove (3)$\Rightarrow$(2): once $M$ is
finitely generated it fits into a sequence as above whose kernel is again
quasiseparated of the form (3), so it too is finitely generated, and $M$ is
finitely presented. To see (3)$\Rightarrow$(2), resolve the tower: choose a
finite free module surjecting onto $M_{1}$, a finite free $R^{\oplus d_{2}}$
onto the relevant kernel, and assemble the finite free modules $R^{\oplus
d_{1}}$, $R^{\oplus d_{1}+d_{2}},\dots$ into a compatible system; in the inverse
limit these give a surjection $\prod_{\mathbb{N}}R\twoheadrightarrow M$. (One
can even arrange a topological splitting of the underlying map of spaces, by
choosing compatible sections, if one does not insist on $R$-linearity.)
\end{proof}

\begin{remark}[Module theory over the endomorphism ring]
\label{rmk:11:endomorphism}
Because $\Solid_{R}$ has a compact projective generator, it is equivalent to
the category of modules over the (noncommutative) endomorphism ring of that
generator, so the extension-closedness in \cref{thm:11:coherence}(2) is just
the usual statement from module theory over an associative ring, and one may
quote it directly.
\end{remark}

\begin{remark}[Finitely presented $\neq$ quasiseparated]
\label{rmk:11:fp-not-qsep}
The lemma does not say finitely presented objects are quasiseparated. A
cokernel of a map $\prod_{\mathbb{N}}R\to\prod_{\mathbb{N}}R$ need not be
quasiseparated: over $\mathbb{Z}$ the module $\mathbb{Z}_{p}/\mathbb{Z}$ is
finitely presented but not quasiseparated. A non-quasiseparated finitely
presented object forces $\dim R\ge 1$; over a finite field every finitely
presented solid module is quasiseparated.
\end{remark}

\begin{remark}[Where finite type is used]
\label{rmk:11:where-finite-type}
Finite type over $\mathbb{Z}$ enters \emph{only} to identify
$\prod_{\mathbb{N}}R$ as the compact projective generator (via the surjection
from a polynomial ring). If one instead \emph{defines} the category by
declaring $\prod_{\mathbb{N}}R$ to be a compact projective generator with the
prescribed endomorphisms, then the rest of \cref{thm:11:coherence,lem:11:qsep-pro} goes through for any noetherian ring.
\end{remark}

\subsection{Flatness of the generator}

\begin{proposition}[$\prod_{\mathbb{N}}R$ is flat]
\label{prop:11:flat}
Let $R$ be finite type over $\mathbb{Z}$ and $M\in\Solid_{R}^{\mathrm{fp}}$.
Then
\begin{equation}
\label{eq:11:flat-formula}
M\otimes^{\mathbb{L}\solid}_{R}\textstyle\prod_{\mathbb{N}}R\ \cong\
\textstyle\prod_{\mathbb{N}}M,
\end{equation}
which in particular lives in degree $0$. Consequently $\prod_{\mathbb{N}}R$ is
flat for the solid tensor product $\otimes^{\solid}_{R}$.
\end{proposition}

\begin{proof}
Flatness follows from \eqref{eq:11:flat-formula}: the derived tensor product
with any finitely presented object lives in degree $0$, and every object is a
filtered colimit of finitely presented ones, so the derived tensor product
with $\prod_{\mathbb{N}}R$ lands in degree $0$ throughout.

For \eqref{eq:11:flat-formula}, resolve $M$ by copies of $\prod_{\mathbb{N}}R$
(an infinite resolution: each successive kernel is quasiseparated, hence
finitely presented by \cref{lem:11:qsep-pro}). This reduces the claim to
$M=\prod_{\mathbb{N}}R$, and then, writing $\prod_{\mathbb{N}}R=(\prod_{\mathbb
{N}}\mathbb{Z}\otimes_{\mathbb{Z},\solid}R)$ solidified, to the analogous fact
for $\prod_{\mathbb{N}}\mathbb{Z}$ over $\mathbb{Z}$: that
$\prod_{\mathbb{N}}\mathbb{Z}$ distributes over infinite products. This was
proved via the light condensed set $P=\mathbb{N}\cup\{\infty\}$, whose self
product $P\times P$ is isomorphic to $P$ in the expected way
(\cref{cor:5:tensor-square,cor:6:flat}).
\end{proof}

\begin{remark}[Defining the derived tensor product]
\label{rmk:11:derived-tensor-def}
Making sense of $\otimes^{\mathbb{L}}$ does not require producing flat objects
by hand. On $D(\CondAbl)$ one derived-sheafifies the section-wise tensor
product of complexes of abelian groups; over a condensed ring one takes the
relative tensor product (a bar construction of derived tensor products), and on
a localization one has the symmetric monoidal structure supplied by the
analytic ring structure. So the construction is available before any flatness
is known. That said, in light condensed abelian groups there \emph{are} enough
flats --- the free object on a light condensed set is flat --- and, once the
flatness and degree-$0$ statements above are in hand, $\otimes^{\mathbb{L}
\solid}_{R}$ is genuinely the left derived functor of $\otimes^{\solid}_{R}$
(resolvable in either variable), and derived solidification is the left derived
functor of solidification.
\end{remark}

\subsection{Homological dimension over a regular ring}

Over $\mathbb{Z}$, Scholze showed that finitely presented solid abelian groups
have a two-term resolution by products of copies of $\mathbb{Z}$
(\cref{rmk:6:length-one}), so solid $\mathbb{Z}$ has homological dimension $1$
on finitely presented objects, exactly as $\mathbb{Z}$ does. Here is the
generalization.

\begin{theorem}[Ext-vanishing in the regular case]
\label{thm:11:regular-ext}
Let $R$ be finite type over $\mathbb{Z}$ and regular of dimension $d$. Then for
all $M\in\Solid_{R}^{\mathrm{fp}}$ and all $N\in\Solid_{R}$ one has
\begin{equation}
\label{eq:11:regular-ext}
\Ext^{i}_{R}(M,N)=0\qquad\text{and even}\qquad \iExt^{i}_{R}(M,N)=0,\qquad
\forall\,i>d.
\end{equation}
\end{theorem}

\begin{corollary}[Tor bound]
\label{cor:11:tor-bound}
For $M,N\in\Solid_{R}$ $(R$ finite type over $\mathbb{Z}$, regular of dimension
$d)$,
\begin{equation}
\label{eq:11:tor-bound}
H_{i}\bigl(M\otimes^{\mathbb{L}\solid}_{R}N\bigr)=0\qquad\forall\,i>d.
\end{equation}
\end{corollary}

Indeed \cref{thm:11:regular-ext} yields, for finitely presented $M$, a
projective resolution of length $\le d$ (or $d+1$, depending on the
convention), and the projective objects are flat (\cref{prop:11:flat}); a
filtered-colimit argument passes from finitely presented to arbitrary $M$.

\begin{remark}[The hypothesis ``finitely presented'' is essential]
\label{rmk:11:not-fp}
\Cref{thm:11:regular-ext} does \emph{not} extend to arbitrary $M$. Unlike the
classical situation over a regular noetherian ring of finite dimension --- where
the $\Ext$-vanishing holds for all modules --- the passage from finitely
presented to general modules requires analysing a derived inverse limit along
an arbitrary filtered system, and that step fails here. As a concrete (and
cute) exercise over $R=\mathbb{Z}$, take two distinct primes $p\ne\ell$ and the
non-finitely-presented module $\mathbb{Z}_{p}/\mathbb{Z}[1/\ell]$; then
\begin{equation}
\label{eq:11:cute-ext}
\Ext^{2}\bigl(\mathbb{Z}_{p}/\mathbb{Z}[1/\ell],\ \mathbb{Z}\bigr)\ =\
\mathbb{Z}_{\ell}/\mathbb{Z}[1/p]\ \neq\ 0,
\end{equation}
symmetric under interchanging the roles of $p$ and $\ell$. Clausen expects
arbitrarily high non-vanishing $\Ext$-groups to occur here even in ZFC (and
they certainly do in some model of ZFC --- the relevant $\Ext$-group is one of
those studied by set theorists, cf.\ \cref{def:4:star,thm:4:star-set-theory}).
If the continuum $2^{\aleph_{0}}$ is bounded by some $\aleph_{n}$, one does
get cohomological-dimension bounds; if $2^{\aleph_{0}}>\aleph_{n}$ for all $n$,
one should expect arbitrarily high $\Ext$.
\end{remark}

\subsubsection*{Proof of \texorpdfstring{\cref{thm:11:regular-ext}}{the theorem}}

\emph{Step 1: $M$ quasiseparated, $N$ discrete.} Take $N$ a discrete finitely
generated $R$-module and $M$ quasiseparated finitely presented, so
$M=\varprojlim_{n}M_{n}$ with surjective transition maps
(\cref{lem:11:qsep-pro}). We claim
\begin{equation}
\label{eq:11:ext-colimit}
\Ext^{i}_{R^{\solid}}(M,N)\ =\ \varinjlim_{n}\Ext^{i}_{R}(M_{n},N).
\end{equation}
(The internal statement follows from the underlying one by replacing $N$ with
$\Cont(S,N)$, another discrete module.) The content is that $R\Homop$ pulls the
inverse limit out. The Mittag-Leffler resolution
\begin{equation}
\label{eq:11:ML-resolution}
\varprojlim_{n}M_{n}\ \longrightarrow\ \textstyle\prod_{n}M_{n}\
\xrightarrow{\ \mathrm{id}-\mathrm{shift}\ }\ \textstyle\prod_{n}M_{n}
\end{equation}
turns the question into whether $R\Homop$ takes the product $\prod_{n}M_{n}$ to
the direct sum; resolving $\prod_{n}M_{n}$ by copies of $\prod_{\mathbb{N}}R$
reduces this to
\begin{equation}
\label{eq:11:hom-product}
R\Homop\bigl(\textstyle\prod_{\mathbb{N}}R,\,N\bigr)\ =\ \bigoplus_{\mathbb{N}}
R\Homop_{R}(R,N),
\end{equation}
which holds because $\prod_{\mathbb{N}}R$ is projective and by the basic
computation of $\Homop$. Finally, $R$ being regular of dimension $d$, the
classical $\Ext$-groups $\Ext^{i}_{R}(M_{n},N)$ between discrete modules vanish
for $i>d$, so the colimit \eqref{eq:11:ext-colimit} vanishes for $i>d$. Thus
the quasiseparated question has been reduced to the classical one in discrete
commutative algebra --- and it is exactly here that regularity is used.

\emph{Step 2: general finitely presented $M$, the naive bound.} For
$M\in\Solid_{R}^{\mathrm{fp}}$ arbitrary, resolve
\begin{equation}
\label{eq:11:naive-resolution}
0\to K\to\textstyle\prod_{\mathbb{N}}R\to M\to 0,
\end{equation}
with $K$ quasiseparated and finitely presented. Chasing the long exact
sequence and using Step~1 only gives
\begin{equation}
\label{eq:11:naive-bound}
\Ext^{i}_{R}(M,N)=0\qquad\forall\,i>d+1,
\end{equation}
one degree worse than wanted. To recover the sharp bound, continue the
resolution one further step (assume $d\ge 2$; the cases $d\le 1$ are handled by
variants of the same argument):
\begin{equation}
\label{eq:11:two-step-resolution}
0\to N\to\textstyle\prod F_{n}\to\textstyle\prod F'_{n}\to M\to 0,
\qquad F_{n},F'_{n}\ \text{finite free},
\end{equation}
so that $N$ is a second syzygy. To prove $\Ext^{i}(M,X)=0$ for $i>d$ it now
suffices to prove
\begin{equation}
\label{eq:11:suffices-syzygy}
\Ext^{i}(N,X)=0\qquad\forall\,i>d-2.
\end{equation}

\emph{Step 3: two presentations of the syzygy.} There are two ways to write the
quasiseparated module $N$ as an inverse limit of finitely generated
$R$-modules, each good in one respect and bad in the other:
\begin{enumerate}
\item $N_{n}:=\ker(F_{n}\to F'_{n})$, so $N=\varprojlim_{n}N_{n}$. \emph{Good:}
each $N_{n}$ is a kernel of a map between finite frees, i.e.\ a second syzygy,
so $\Ext^{i}(N_{n},X)=0$ for $i>d-2$. \emph{Bad:} no guarantee that $(N_{n})$
is Mittag-Leffler --- which is what \eqref{eq:11:ext-colimit} needs to compute
$\Ext^{i}(N,X)=\varinjlim_{n}\Ext^{i}(N_{n},X)$.
\item $N'_{n}:=\operatorname{im}\bigl(N\to F_{n}\bigr)\subseteq N_{n}$.
\emph{Good:} Mittag-Leffler, indeed surjective transition maps by construction.
\emph{Bad:} no guarantee that $\Ext^{d-1}(N'_{n},X)=0$.
\end{enumerate}

\emph{Step 4: interpolating.} One finds $N'_{n}\subseteq N''_{n}\subseteq N_{n}$
with both good properties:
\begin{equation}
\label{eq:11:interpolate}
\Ext^{i}(N''_{n},X)=0\ \ \forall\,i\ge d-1,\qquad\text{and}\qquad (N''_{n})\
\text{is Mittag-Leffler.}
\end{equation}
Then $\varprojlim_{n}N''_{n}=N$ (it is sandwiched between two systems with the
same inverse limit), and \eqref{eq:11:interpolate} yields
\eqref{eq:11:suffices-syzygy}.

The obstruction to $N'_{n}$ itself having $\Ext^{d-1}(N'_{n},X)=0$ for all $X$
is a matter of \emph{depth}. By the Auslander--Buchsbaum formula, for a regular
noetherian local ring
\begin{equation}
\label{eq:11:auslander-buchsbaum}
\pdim(M)+\depth(M)=\dim,
\end{equation}
and because we already have the one-better estimate on projective dimension,
the only obstruction is at the maximal ideals whose local ring has the full
dimension $d$, where the issue is passing from depth $1$ to depth $2$: at such
a maximal ideal $N'_{n}$ is only known to have depth $1$.

To fix this, suppose (for simplicity) the problem occurs at a single closed
point $x\in\Spec R$; in general one fixes finitely many, and by
noetherianness the situation stabilizes. With $j\colon\Spec R\setminus\{x\}
\hookrightarrow\Spec R$, replace $N'_{n}$ by $j_{*}j^{*}N'_{n}$. There is a map
$N'_{n}\to j_{*}j^{*}N'_{n}$, which is an inclusion by the depth-$1$ hypothesis;
and, since local cohomology at the maximal ideal computes depth, this operation
improves the depth at $x$ to $2$. Performed on $N_{n}$ (already of depth $2$) it
would change nothing, so the result $N''_{n}$ genuinely sits between $N'_{n}$
and $N_{n}$, and fixes the problem at $x$. Doing this compatibly at all the
finitely many bad closed points produces the tower $(N''_{n})$ with depth $2$
everywhere --- hence \eqref{eq:11:interpolate}'s $\Ext$-vanishing.

For the Mittag-Leffler property of $(N''_{n})$, use the short exact sequence of
towers
\begin{equation}
\label{eq:11:ML-ses}
0\to (N'_{n})\to (N''_{n})\to (N''_{n}/N'_{n})\to 0.
\end{equation}
The subtower $(N'_{n})$ is Mittag-Leffler (surjective transitions), and the
quotient tower $(N''_{n}/N'_{n})$ is supported at the finitely many closed
points and finitely generated over $R$, hence \emph{finite} as an abelian group
(the residue fields are finite, $R$ being finite type over $\mathbb{Z}$); an
inverse system of finite abelian groups is automatically Mittag-Leffler by a
compactness argument. Thus $(N''_{n})$ is Mittag-Leffler.

\emph{Step 5: bootstrapping the coefficients.} With Steps 1--4 the theorem holds
whenever $M$ is finitely presented and $N=X$ is a discrete module. Now let $X$
be quasiseparated finitely presented, $X=\varprojlim_{n}X_{n}$ with surjective
transitions; then $R\Homop(M,X)=R\varprojlim_{n}R\Homop(M,X_{n})$. A priori the
$\varprojlim$ could push one degree higher, but the possible $\varprojlim^{1}$
lives on the $\Ext^{d}$-groups, and there the surjectivity of the transition
maps makes the maps on $\Ext^{d}$ surjective (the obstruction is an
$\Ext^{d+1}$ into a kernel, which vanishes), so no $\varprojlim^{1}$ term
contributes in degree $d+1$. Hence the bound $i>d$ persists for such $X$. For
$X$ arbitrary finitely presented one resolves by quasiseparated modules; for
$X$ a filtered colimit $\varinjlim_{j}X_{j}$ of finitely presented modules one
uses pseudocoherence of $M$ (finitely presented objects have infinite
resolutions by copies of $\prod_{\mathbb{N}}R$, which are compact projective) to
commute $\Ext^{i}(M,-)$ with the colimit. Finally, the internal $\iExt^{i}$
statements follow from the underlying ones by passing to $S$-valued points, and
the $\varprojlim^{1}$ arguments go through at the internal level because in
condensed abelian groups the condensed $\varprojlim^{1}$ is the sheafification
of the section-wise $\varprojlim^{1}$, which can be computed termwise since
countable products are exact. \qed

\subsection{A second localization: the constructible topology}

We turn from structure theory to geometry. Recall from \cref{lec:9} that for a
discrete commutative ring $R$ the derived category $D((R,\mathbb{Z})^{\solid})$
localizes over the valuative spectrum $\Spv(R)$. That localization is convenient
because its unit, restricted to a basic quasicompact open, stays discrete: over
a rational open $U(f_{1},\dots,f_{n}/g)$ the structure ring is simply the
discrete localization $R[\tfrac1g]$. It keeps one close to Huber's world. But it
is by no means the only localization of this category, and we now describe a
genuinely different one, departing more radically from Huber's formalism.

\begin{definition}[Idempotent algebra of a basic constructible subset]
\label{def:11:constructible-idempotent}
Let $C\subseteq\Spec R$ be a constructible subset, that is, a finite union of
locally closed subsets, where a basic locally closed subset is an intersection
\begin{equation}
\label{eq:11:basic-constructible}
D(g)\cap Z(f_{1},\dots,f_{n})\ =\ \{g\neq 0\}\cap\{f_{1}=\cdots=f_{n}=0\}
\end{equation}
of a distinguished quasicompact open with the common zero locus of finitely many
functions. To such a basic locally closed subset one assigns the object
\begin{equation}
\label{eq:11:constructible-algebra}
A\ =\ R\bigl[\tfrac1g\bigr]^{\wedge}_{(f_{1},\dots,f_{n})}\ \in\ D(\Solid_{
\mathbb{Z}}),
\end{equation}
the $(f_{1},\dots,f_{n})$-\emph{derived} completion of $R[\tfrac1g]$, formed in
the derived category of solid abelian groups --- i.e.\ one inverts $g$ and then
puts the inverse-limit (adic) topology on the derived inverse limit. This object
depends only on the constructible subset. When $R$ is noetherian the derived
completion agrees with the ordinary completion.
\end{definition}

\begin{remark}[Only locally closed pieces are idempotent]
\label{rmk:11:only-locally-closed}
One should attach the idempotent algebra only to a basic locally closed subset,
not to an arbitrary constructible one. For a general constructible $C$, written
as a union of basic locally closed pieces, one takes the corresponding limit of
the pieces (and their intersections, which are tensor products,
\cref{rmk:11:idempotent-tensor}); the resulting object $R\Gamma$ can have higher
cohomological degree, depending on the number of pieces used to cover $C$. The
only derived behaviour arises from a union of principal opens.
\end{remark}

\begin{remark}[Idempotence]
\label{rmk:11:idempotent-tensor}
Recall (\cref{prop:6:tensor-p-complete} and its generalization) that the solid
tensor product of \emph{connective} derived complete objects is derived complete. Writing
$A=R[\tfrac1g]^{\wedge}_{(f_{1},\dots,f_{n})}$, idempotence
\begin{equation}
\label{eq:11:idempotent}
A\otimes^{\mathbb{L}\solid}_{(R,\mathbb{Z})^{\solid}}A\ \xrightarrow{\ \sim\ }\ A
\end{equation}
is checked modulo $(f_{1},\dots,f_{n})$: since both sides are derived complete
one may reduce modulo the sequence, whereupon the statement becomes the obvious
fact that a localization of a ring, tensored over the ring with itself, is that
localization again. Intersections of basic constructible subsets correspond to
solid tensor products of the associated algebras.
\end{remark}

\begin{proposition}[Constructible covers give a descent datum]
\label{prop:11:constructible-cover}
Let $C_{1},\dots,C_{n}\subseteq\Spec R$ be locally closed constructible subsets
with $\bigcup_{i}C_{i}=\Spec R$, and $A_{1},\dots,A_{n}$ the associated
idempotent algebras. Then the $A_{i}$ \emph{cover} $D((R,\mathbb{Z})^{\solid})$:
\begin{equation}
\label{eq:11:cover-conservative}
M\in D\bigl((R,\mathbb{Z})^{\solid}\bigr),\quad M\otimes A_{i}=0\ \forall i\
\Longrightarrow\ M=0,
\end{equation}
and consequently $D((R,\mathbb{Z})^{\solid})$ is recovered as the totalization
of the associated \v{C}ech diagram
\begin{equation}
\label{eq:11:cech}
D\bigl((R,\mathbb{Z})^{\solid}\bigr)\ =\ \varprojlim\Bigl(\ \textstyle\prod_{i}
\Mod_{A_{i}}\ \rightrightarrows\ \prod_{i<j}\Mod_{A_{i}\otimes A_{j}}\
\Rrightarrow\ \cdots\ \Bigr).
\end{equation}
\end{proposition}

\begin{proof}[Idea]
It suffices to show that $R$ is generated, under the triangulated operations, by
$A_{i}$-modules for varying $i$. Consider the simplest cover $\Spec R=D(f)\cup
Z(f)$. Here $R[\tfrac1f]$ is already an $A_{D(f)}$-module, while the fibre of
$R\to R[\tfrac1f]$ is built from $R/f^{n}$'s and is an
$R^{\wedge}_{f}$-module, hence lies over $Z(f)$. So $R$ is generated by an
$A_{D(f)}$-module and an $A_{Z(f)}$-module; tensoring, one hits every object,
and \eqref{eq:11:cover-conservative} follows.
\end{proof}

\begin{example}[Higher local fields]
\label{ex:11:higher-local-fields}
Take $R=\mathbb{Z}[X,Y]$ and cover $\Spec R$ by three locally closed pieces
--- successively imposing $X\neq0$; then $X=0,\ Y\neq0$; then $X=Y=0$:
\begin{align}
\label{eq:11:hlf-pieces}
C_{1}&=\{X\neq 0\}, & A_{1}&=\mathbb{Z}[X^{\pm1},Y],\notag\\
C_{2}&=\{X=0,\ Y\neq 0\}, & A_{2}&=\mathbb{Z}[Y^{\pm1}][\![X]\!],\\
C_{3}&=\{X=Y=0\}, & A_{3}&=\mathbb{Z}[\![X,Y]\!].\notag
\end{align}
The interesting object is the triple intersection. Tensoring $A_{2}$ with
$A_{1}$ inverts $X$, giving $\mathbb{Z}[Y^{\pm1}](\!(X)\!)$; tensoring further
over $\mathbb{Z}[X,Y]$ (in $\Solid_{\mathbb{Z}}$) with $A_{3}=\mathbb{Z}[\![X,Y]
\!]$ inverts $Y$ as well, and one computes
\begin{equation}
\label{eq:11:hlf-triple}
A_{C_{1}\cap C_{2}\cap C_{3}}\ =\ \mathbb{Z}(\!(Y)\!)(\!(X)\!),
\end{equation}
each such value obtained by iterating ``complete, then invert''. The order in
which one completes and inverts famously matters, yet here it all arises from a
manifestly commutative construction. These objects are the higher local fields
that have been studied classically; as topological rings or fields they are
notoriously ill-behaved, but in the condensed world they are perfectly
well-behaved objects, arising from this natural, idempotent (even derived) and
functorial procedure.%
\footnote{Clausen attributes to A.~Mathew a paper documenting how badly higher
local fields behave as topological rings; the precise reference was not given.}
\end{example}

\begin{remark}[Where the solid $p$-adic functional analysis is written up]
\label{rmk:11:references}
The solid functional analysis underlying the ``derived complete'' facts used
above (\cref{rmk:11:idempotent-tensor}) --- Scholze's treatment from earlier in
the course --- is written down by Bosco, in an appendix on $p$-adic functional
analysis from the solid perspective \cite{bosco2023rationalpadichodgetheory}
(the $p$-complete case), and in general by Mann in his thesis
\cite{mann2022padic6functorformalismrigidanalytic}.
\end{remark}
